\section{The methods of version 2.0}\label{sec:methods}

All methods of this section are those of the PhD dissertation, written in a common form. The
FDS-controlled methods (\code{gaussian}, \code{kfds}, \code{rf}, \code{pf}, \code{psf}, \code{sor})
end with the closed loop of Section~\ref{sec:closedloop}.

\subsection{kFDS: prescribed FDS and kurtosis}\label{sec:kfds}

\emph{Form of the signal.} kFDS multiplies a stationary Gaussian carrier $g$, designed as in
Section~\ref{sec:inversion}, by a slowly varying envelope $e(t)$,~\eqref{eq:envcarrier}. The envelope
is piecewise constant over windows of length $L_w$, with linear ramps of $L_w/3$ between the levels, so
that the waveform has no jumps.

\emph{The damage is kept by the normalisation $\tavg{e^b} = 1$.} If a window is long compared with the
response time of the oscillators, the response in the window is the carrier's response multiplied by
$e_w$, its cycles are those of the carrier scaled by $e_w$, and by~\eqref{eq:minerz} the damage in the
window is $e_w^b$ times the carrier's damage in the window. If the carrier's damage is spread evenly
over time, summing over the windows gives
\begin{equation}
  D(e\,g) = \sum_w e_w^b\,D_{g,w} = \sum_w e_w^b\,\frac{T_w}{T}D_g = \tavg{e^b}\,D_g ,
\end{equation}
so any envelope with $\tavg{e^b}=1$ keeps the FDS of the carrier. (In the dissertation the same
condition appears as weights $v_j$ of blocks, $e_j = (n v_j)^{1/b}$.)

\emph{The kurtosis is set by the contrast of the envelope.} By~\eqref{eq:kmod} the kurtosis of the
product is $3\tavg{e^4}/\tavg{e^2}^2$. The envelope is $e_w = \exp(\alpha u_w)$ normalised to
$\tavg{e^b}=1$, where the relative log-levels $u_w\le0$ are drawn
\begin{itemize}[nosep]
\item from the distribution of the window RMS of the reference (\code{pattern = "reference"},
default when a record is available): the logarithm of the RMS of every window of the record is
computed, and $n_w$ values are drawn from their life-weighted distribution by quantile resampling, in
random order, so that the test has bursts with the statistics of the field data;
\item or as random bursts (\code{"bursts"}): 10 to 20\,\% of the windows at a high level.
\end{itemize}
$\alpha = 0$ gives a constant envelope (kurtosis 3); the kurtosis grows with $\alpha$. The program
finds the $\alpha$ that gives the target kurtosis \emph{on the realised signal} $e\,g$ with a
bracketing root finder (Brent's method). The envelope is frozen after the first two iterations of the
closed loop, which then corrects only the carrier spectrum.

\emph{Window length.} The normalisation argument requires windows long compared with the response
time of the slowest oscillator. The amplitude of a free oscillator decays as $e^{-\zeta\wn t}$, with the
decay time $1/(\zeta\wn) = Q/(\pi\fn) = 2\tau$; the window is at least $2Q/(\pi f_{\min})$. Within this
bound the default follows the rule of the dissertation: a fraction of the reference duration (1/150 to
1/400 depending on $f_s$), with at least 300 samples.

\subsection{RF: prescribed FDS and RMS}

RF uses the same envelope with the RMS as the target. With $\tavg{e^b}=1$ the RMS of the product is
$\sigma_g\tavg{e^2}^{1/2}$. For $b>2$ the power-mean inequality gives $\tavg{e^2}^{1/2}\le
\tavg{e^b}^{1/b} = 1$, with equality only for a constant envelope; increasing the contrast lowers the
RMS. RF can therefore only lower the RMS below that of the Gaussian profile with the same FDS; the
program says so when a higher RMS is requested.

\subsection{PF and PSF: prescribed FDS and PSD}\label{sec:pf}

PF builds the test from two independent stationary Gaussian carriers, a high-energy one $x_B$ (PSD
$G_B$) for the bursts and a low-energy one $x_L$ (PSD $G_L$) for the rest, combined with a smooth
switching mask $m(t)\in[0,1]$ that is 1 during the bursts:
\begin{equation}
  x(t) = \sqrt{m(t)}\,x_B(t) + \sqrt{1-m(t)}\,x_L(t).
\end{equation}
Because the carriers are independent, the local PSD is $mG_B + (1-m)G_L$, and the time-averaged PSD is
$\tavg{m}G_B + (1-\tavg m)G_L = (t_BG_B + t_LG_L)/T$, with $t_B = \tavg mT$ the total burst time and
$t_L = T-t_B$. Requiring the reference PSD gives
\begin{equation}
  G_L = \frac{T\,G_{\mathrm{ref}} - t_B\,G_B}{t_L}.
\end{equation}
The damage is the sum of the damages of the two parts. The burst PSD is found by a fixed point:
start with $D_L = D_t/20$; design $G_B$ for $D_B = D_t-D_L$ over $t_B$ (Section~\ref{sec:inversion});
compute $G_L$ from the PSD condition and its Gaussian damage $D_L$ over $t_L$; repeat. The burst time
starts at $T/20$ and is halved while $G_L$ would be negative over more than 10\,\% of the band, but
two bursts of at least three decay times of the lowest oscillator are kept when possible. When
PSD, RMS and FDS cannot all be met (for instance at high acceleration), the program reports it.

PSF prescribes only the shape of the PSD: PF is run with the damage that the reference produces in
the test duration, $D_t\,T/\Tl$, and the result is multiplied by $(\Tl/T)^{1/b}$, which multiplies
its damage by $\Tl/T$ and its PSD by $(\Tl/T)^{2/b}$ without changing the shape.

\subsection{Kurtosis control without damage: MLV, VSD, PM}

These methods prescribe the PSD and the kurtosis of the signal, not its damage. MLV
(multi-level variance) is the modulated Gaussian~\eqref{eq:envcarrier} normalised to $\tavg{e^2}=1$
(this keeps PSD and RMS) with the contrast solved for the kurtosis. VSD (variable spectral density)
splits the band into log-spaced sub-bands, each with its own carrier and independent burst envelope,
so that the shape of the PSD changes with time while its time average is the reference PSD. PM
(phase manipulation) keeps Gaussian amplitudes and partially aligns the phases of the harmonics in a
fraction of the blocks, which creates peaks; the degree of alignment is solved for the kurtosis. The
benchmark of the program shows that, over the same duration, their FDS can differ from the
reference's by a factor of 3 or more: kurtosis alone does not control the damage.

\subsection{MRS control}

The largest value of a narrow-band Gaussian response in the duration $T$ can be estimated from Rice's
crossing rate. Repeating the derivation of~\eqref{eq:nu0} for the level $z=u$ instead of 0, the
joint density at $z=u$ contains the extra factor $e^{-u^2/(2\sigma_z^2)}$, so the rate of up-crossings
of $u$ is $\nu_0^+e^{-u^2/(2\sigma_z^2)}$. The level crossed on average once in $T$ satisfies
$\nu_0^+Te^{-u^2/(2\sigma_z^2)} = 1$, i.e.
\begin{equation}
  u = \sigma_z\sqrt{2\ln(\nu_0^+T)} \approx \sigma_z\sqrt{2\ln(\fn T)} ,
  \qquad \mathrm{MRS}_G(\fn) = \wn^2\sigma_z\sqrt{2\ln(\fn T)} .
\end{equation}
A target MRS therefore prescribes $\sigma_z$ at every node, the PSD follows from~\eqref{eq:WG}, and
the closed loop corrects the measured MRS with the exponent 1 instead of $1/b$ in~\eqref{eq:gain}
(the MRS is proportional to the amplitude).

\subsection{Sine on random}\label{sec:sor2}

Rotating machines add tonal lines to the random vibration. The method \code{sor} detects the lines
in the reference PSD: a peak standing more than a prominence threshold (10\,dB by default) above a
running median of the PSD is a line, and its amplitude follows from the power in the peak,
$A = \sqrt{2P}$ (a sine of amplitude $A$ has the power $A^2/2$). A stationary sine of amplitude $A$
produces in an oscillator a sinusoidal response of amplitude $|H|A$ at the line frequency, i.e.\
the damage $\frac{K^b}{C}f_0 T(|H|A)^b$ in the time $T$. To keep the damage of the life $\Tl$ in the
test duration $T$, the test amplitude $A_T$ must satisfy $A_T^bT = A^b\Tl$:
\begin{equation}\label{eq:sorgain}
  A_T = A\left(\frac{\Tl}{T}\right)^{1/b}.
\end{equation}
The random part is designed for the remaining damage and corrected in closed loop with the lines kept
fixed; nodes where the lines alone exceed the target are reported and left out of the loop criterion.
Version~2.1 adds the detection of lines that are present only in some blocks of a record
(Section~\ref{sec:lines}).

\subsection{Missions}\label{sec:missions}

A mission is a tree whose leaves are situations (a record or a specification with the life it
represents) and whose inner nodes are groups. In a \emph{series} group the component lives through
all children one after the other: the FDS are added, the lives are added, the PSD is the
life-weighted mean, the MRS is the maximum. In a \emph{parallel} group it lives through only one of
them, not known in advance: the FDS and PSD are enveloped (maximum over the children) and the life,
RMS and kurtosis are those of the most damaging child.

The statistics of a series group follow from the moments of a mixture. If a fraction $w_c$ of the
life is spent in situation $c$, where the signal has the moments $\mu_{2,c}$ and $\mu_{4,c}$, the
moments over the whole life are the weighted averages $\mu_2 = \sum_cw_c\mu_{2,c}$ and $\mu_4 =
\sum_cw_c\mu_{4,c}$, and the kurtosis is $\mu_4/\mu_2^2$. Version~2.1 applies the same rule to the
response of every oscillator: with $v_c(\fn)$ and $\kappa_c(\fn)$ the variance and kurtosis of the
response in situation $c$, $\mu_{4,c} = \kappa_cv_c^2$ and
\begin{result}[Response kurtosis of a series mission]
\begin{equation}\label{eq:kzmix}
  v(\fn) = \sum_c w_c v_c(\fn),\qquad
  \kappa_z(\fn) = \frac{\sum_c w_c\,\kappa_c(\fn)\,v_c(\fn)^2}{v(\fn)^2}.
\end{equation}
\end{result}
A PSD specification has $\kappa_c = 3$ (it describes a Gaussian process); a swept sine is measured on
its generated signal.
